\chapter[Governance, Reproducibility, and Conclusion]{Governance, Reproducibility, Limitations, and Conclusion}
\label{ch:governance-reproducibility-limitations-conclusion}

\section{What is proved}
The dissertation proves a bounded but durable theory claim. Under explicit assumptions,
degraded observation can create a gap between observed-state legality and true-state safety;
reliability-conditioned uncertainty inflation changes the admissible release set; and
repair-or-fallback discipline preserves safety whenever the inflated set and lifecycle
semantics remain feasible. Those results live in the theorem chapter and proof appendix and
should not be overstated beyond their assumptions.

\section{What is empirically validated}
The empirical surface is strongest at the witness and defended bounded rows. Battery remains
the deepest theorem-to-code-to-artifact chain. Industrial process control and healthcare
monitoring support defended bounded interpretations under the current replay and governance
surfaces. The autonomous-vehicle row is a defended bounded runtime row under the present
artifact surface. Navigation remains blocked by the absence of a defended real-data
train--replay chain. Aerospace remains blocked by the absence of a defended multi-flight
runtime validation surface with material post-repair gain.

Those limits matter because ORIUS is not presented here as a universal optimal controller or
as a finished equal-domain deployment program. It is presented as a runtime safety layer
whose strongest current evidence is deep in one witness domain and bounded across several
others. The open rows should therefore be read as the next empirical work required to extend
the thesis, not as defects that invalidate the defended core.

\section{What is out of scope}
The dissertation does not claim:
\begin{itemize}
  \item equal-domain empirical closure across all six domains;
  \item universal statistical calibration under arbitrary shift;
  \item production deployment readiness in every plant family; or
  \item formal certification compliance under UL 4600, NIST AI RMF, ISO 26262, or IEC 61508.
\end{itemize}
The standards language used in the governance chapter is an evidence-structure alignment and
lifecycle-analogy argument only. The deployment language is deliberately bounded by the
parity matrix, runtime/governance surfaces, and explicit non-claim statements.

\section{Reproducibility discipline}
Every headline claim stays tied to tracked figures, tables, manifests, or claim matrices so
that the defended surface remains auditable after the fact. The purpose of that discipline is
simple: a universal safety thesis should be inspectable at the level of artifacts, not only
persuasive at the level of prose.

That is why the dissertation remains chapter-first in prose but artifact-first in evidence.
The prose chapters explain the argument. The claim validator, publication matrices, release
manifests, and replay tables control what the prose is allowed to say. In that exact sense,
the dissertation is universal in architecture, tiered in evidence.

\section{Threats to validity}
The largest threats to validity are asymmetry and interpretation. Asymmetry arises because
the witness row is much deeper than the currently blocked rows. Interpretation risk arises
because readers may confuse structural universality with equal-domain parity. The
dissertation addresses the first risk by keeping the witness role explicit and by preserving
the parity matrix in the main narrative. It addresses the second risk by separating proved,
validated, and roadmap claims at the major boundary chapters.

\section{Conclusion}
The argument of this dissertation can now be stated plainly. Degraded observation is not a
peripheral nuisance attached to otherwise reliable autonomy; it is a release-boundary safety
problem in its own right. Once observed-state legality and true-state legality can diverge,
safe actuation requires more than nominal intelligence. It requires explicit reliability
estimation, conservative state widening, repaired action or bounded fallback, and
certificate-aware release.

ORIUS makes that claim operational through three connected surfaces. A
\texttt{DomainAdapter} exposes the local safety predicate and fallback object. The universal
benchmark schema exposes one replay contract for truth, observation, intervention, fallback,
certificate validity, and latency. CertOS supplies the lifecycle semantics that determine
when repaired action may still be released. Together, those surfaces turn degraded
observation from a vague systems concern into a measurable and governable runtime object.

That is the durable contribution of the thesis. ORIUS does not ask every domain to share a
controller, plant model, or optimization objective. It asks every domain to make the
release boundary explicit. In doing so, it turns true-state violation, intervention burden,
fallback usage, certificate validity, and runtime latency into scientific objects that can
be compared across domains without pretending the domains themselves are identical.

The final thesis claim is therefore ambitious and disciplined at the same time. ORIUS
already supports a reusable runtime safety layer for physical AI under degraded observation.
It does not yet support the stronger claim that every supported domain is equally mature. In
that exact sense, structural universality is defended now, while equal-domain closure remains an empirical program that must be earned. That boundary is not a rhetorical
retreat. It is part of what makes the argument credible.
